\documentclass[a4paper,11pt,reqno]{amsart}
\usepackage[utf8]{inputenc}
\usepackage{enumerate}
\usepackage{amsmath,amssymb,amsthm}
\usepackage{mathtools}
\usepackage{xcolor}
\usepackage{hyperref}

\newtheorem{lemma}{Lemma}[section]
\newtheorem{theorem}{Theorem}
\newtheorem{corollary}{Corollary}[section]
\newtheorem{proposition}{Proposition}[section]
\newtheorem{remark}{Remark}[section]

\DeclareMathOperator{\C}{\mathbb{C}}
\DeclareMathOperator{\R}{\mathbb{R}}
\newcommand{\N}{\mathbb{N}}
\newcommand{\Z}{\mathbb{Z}}

\begin{document}

\begin{lemma}[Convex sets of fixed area with arbitrarily small localization norm]
Fix \(m>0\). Then there exists a family of bounded convex sets
\(\{R_L\}_{L\geq 1}\subset \C\) such that
\[
|R_L|=m
\qquad\text{for every }L\geq 1,
\]
and, for the Gaussian localization operator in the Bargmann--Fock model,
\[
\|\mathcal L_{R_L}\|_{\mathcal F^2\to \mathcal F^2}\to 0
\qquad\text{as }L\to\infty.
\]
More precisely, one may take the rectangles
\[
R_L:=
\left[-\frac L2,\frac L2\right]\times
\left[-\frac{m}{2L},\frac{m}{2L}\right],
\]
for which
\[
\|\mathcal L_{R_L}\|_{\mathcal F^2\to \mathcal F^2}^2
\leq
\|\mathcal L_{R_L}\|_{HS}^2
\lesssim
\frac{m^2}{L}.
\]
\end{lemma}

\begin{proof}
We work in the Bargmann--Fock representation, where
\[
(\mathcal L_\Omega F)(w)
=
\int_\Omega F(z)e^{\pi \overline zw}e^{-\pi |z|^2}\,dA(z),
\qquad F\in \mathcal F^2(\C).
\]
For such operators one has the Hilbert--Schmidt bound
\[
\|\mathcal L_\Omega\|_{\mathcal F^2\to \mathcal F^2}
\leq
\|\mathcal L_\Omega\|_{HS},
\]
and a direct computation with the reproducing kernel gives
\begin{equation}\label{eq:hs-rectangle}
\|\mathcal L_\Omega\|_{HS}^2
=
\iint_{\Omega\times\Omega} e^{-\pi |z-w|^2}\,dA(z)\,dA(w).
\end{equation}

Now fix \(m>0\), and for \(L\geq 1\) set
\[
R_L:=
\left[-\frac L2,\frac L2\right]\times
\left[-\frac{m}{2L},\frac{m}{2L}\right].
\]
Then each \(R_L\) is bounded and convex, and
\[
|R_L|=L\cdot \frac{m}{L}=m.
\]

Using \eqref{eq:hs-rectangle} and the change of variables \(h=z-w\), one may rewrite
\[
\|\mathcal L_{R_L}\|_{HS}^2
=
\int_{\R^2} e^{-\pi |h|^2}\,|R_L\cap (R_L-h)|\,dh.
\]
Writing \(h=(x,y)\), the overlap of the rectangle with its translate is explicit:
\[
|R_L\cap (R_L-(x,y))|
=
(L-|x|)_+\left(\frac{m}{L}-|y|\right)_+,
\]
where \(t_+:=\max\{t,0\}\). Therefore
\[
\|\mathcal L_{R_L}\|_{HS}^2
=
\int_{-L}^{L}\int_{-m/L}^{m/L}
e^{-\pi(x^2+y^2)}
(L-|x|)\left(\frac{m}{L}-|y|\right)\,dy\,dx.
\]
Since \(e^{-\pi(x^2+y^2)}\le e^{-\pi x^2}\), we obtain
\[
\|\mathcal L_{R_L}\|_{HS}^2
\le
\left(\int_{-L}^{L} (L-|x|)e^{-\pi x^2}\,dx\right)
\left(\int_{-m/L}^{m/L}\left(\frac{m}{L}-|y|\right)\,dy\right).
\]
The second factor is computed exactly:
\[
\int_{-m/L}^{m/L}\left(\frac{m}{L}-|y|\right)\,dy
=
\left(\frac{m}{L}\right)^2.
\]
For the first factor,
\[
\int_{-L}^{L} (L-|x|)e^{-\pi x^2}\,dx
\le
L\int_{\R} e^{-\pi x^2}\,dx
=
L.
\]
Hence
\[
\|\mathcal L_{R_L}\|_{HS}^2
\le
L\left(\frac{m}{L}\right)^2
=
\frac{m^2}{L}.
\]
In particular,
\[
\|\mathcal L_{R_L}\|_{\mathcal F^2\to \mathcal F^2}
\le
\|\mathcal L_{R_L}\|_{HS}
\le
\frac{m}{\sqrt L}
\to 0,
\]
which proves the claim.
\end{proof}

\begin{remark}
This shows that the operator norm of a Gaussian localization operator cannot be
controlled from below in terms of the area of the symbol, even within the class
of bounded convex sets. The same construction also gives simply connected
examples after replacing the rectangle by any sufficiently elongated convex body
of fixed area.
\end{remark}

\begin{proposition}[Simply connected sets far from disks with an eigenfunction close to a prescribed monomial]
Fix \(n\in\N\cup\{0\}\), \(\varepsilon>0\), and \(0<\alpha<2\). Let
\[
e_n(w):=\sqrt{\frac{\pi^n}{n!}}\,w^n
\]
be the normalized monomial in \(\mathcal F^2(\C)\). Then there exists a bounded
simply connected domain \(\Omega\subset\C\) and an \(L^2\)-normalized
eigenfunction \(F_\Omega\) of \(\mathcal L_\Omega\) such that
\[
\|F_\Omega-e_n\|_{\mathcal F^2}<\varepsilon
\]
and
\[
\mathcal A(\Omega)\ge \alpha,
\]
where
\[
\mathcal A(\Omega):=
\inf_{\substack{D\text{ disk}\\ |D|=|\Omega|}}
\frac{|\Omega\triangle D|}{|\Omega|}.
\]
In fact, one may construct a family \(\{\Omega_L\}_{L\ge 1}\) such that
\[
\|F_L-e_n\|_{\mathcal F^2}\to 0
\qquad\text{and}\qquad
\liminf_{L\to\infty}\mathcal A(\Omega_L)\ge \frac{2m}{\pi+m},
\]
for any prescribed \(m>0\).
\end{proposition}

\begin{proof}
Fix \(m>0\), to be chosen later, and let
\[
D:=D_1(0).
\]
Since \(D\) is radial, the localization operator \(\mathcal L_D\) is diagonal in
the orthonormal basis \(\{e_k\}_{k\geq 0}\), and
\[
\mathcal L_D e_k=\lambda_k e_k,
\qquad
\lambda_k=\int_D |e_k(z)|^2e^{-\pi |z|^2}\,dA(z).
\]
Equivalently,
\[
\lambda_k
=
1-e^{-\pi}\sum_{j=0}^k \frac{\pi^j}{j!},
\]
so that
\[
\lambda_k-\lambda_{k+1}
=
e^{-\pi}\frac{\pi^{k+1}}{(k+1)!}>0.
\]
Hence \(\lambda_n\) is a simple isolated eigenvalue of \(\mathcal L_D\).

For \(L\geq 1\), let
\[
R_L:=
\left[L,L+L\right]\times
\left[-\frac{m}{2L},\frac{m}{2L}\right],
\]
let
\[
C_L:=[1,L]\times \left[-\frac{1}{L^2},\frac{1}{L^2}\right],
\]
and define
\[
\Omega_L:=D\cup C_L\cup R_L.
\]
The corridor \(C_L\) meets \(D\) along the segment
\[
\{1\}\times \left[-\frac{1}{L^2},\frac{1}{L^2}\right],
\]
and it also meets \(R_L\). Hence \(\Omega_L\) is bounded, connected, and simply
connected. Moreover,
\[
|R_L|=m,
\qquad
|C_L|=\frac{2(L-1)}{L^2}\le \frac{2}{L}.
\]

Now
\[
\mathcal L_{\Omega_L}
=
\mathcal L_D+\mathcal L_{C_L\setminus D}+\mathcal L_{R_L}.
\]
Since the operator norm is subadditive and \(C_L\setminus D\subset C_L\), we get
\[
\|\mathcal L_{\Omega_L}-\mathcal L_D\|_{\mathcal F^2\to\mathcal F^2}
=
\|\mathcal L_{C_L\setminus D}+\mathcal L_{R_L}\|_{\mathcal F^2\to\mathcal F^2}
\le
\|\mathcal L_{C_L}\|_{\mathcal F^2\to\mathcal F^2}
+
\|\mathcal L_{R_L}\|_{\mathcal F^2\to\mathcal F^2}.
\]
By the elementary bound \(\|\mathcal L_\Omega\|\le |\Omega|\), we have
\[
\|\mathcal L_{C_L}\|_{\mathcal F^2\to\mathcal F^2}\le |C_L|\le \frac{2}{L},
\]
while the rectangle lemma gives
\[
\|\mathcal L_{R_L}\|_{\mathcal F^2\to\mathcal F^2}\le \frac{m}{\sqrt L}.
\]
Therefore
\[
\|\mathcal L_{\Omega_L}-\mathcal L_D\|_{\mathcal F^2\to\mathcal F^2}\to 0.
\]

Let
\[
\gamma_n:=\min_{k\neq n}|\lambda_n-\lambda_k|>0.
\]
For \(L\) sufficiently large,
\[
\|\mathcal L_{\Omega_L}-\mathcal L_D\|<\frac{\gamma_n}{4}.
\]
By standard perturbation theory for isolated simple eigenvalues of compact
self-adjoint operators, there exists then a unique eigenvalue
\(\lambda_{n,L}\) of \(\mathcal L_{\Omega_L}\) in
\((\lambda_n-\gamma_n/2,\lambda_n+\gamma_n/2)\), and the corresponding
rank-one spectral projection \(P_{n,L}\) satisfies
\[
\|P_{n,L}-P_n\|_{\mathcal F^2\to\mathcal F^2}\to 0,
\]
where \(P_nF=\langle F,e_n\rangle e_n\). In particular,
\[
F_L:=\frac{P_{n,L}e_n}{\|P_{n,L}e_n\|_{\mathcal F^2}}
\]
is a well-defined normalized eigenfunction of \(\mathcal L_{\Omega_L}\), and
\[
\|F_L-e_n\|_{\mathcal F^2}\to 0.
\]

It remains to estimate the Fraenkel asymmetry. Set
\[
A_L:=|\Omega_L|=\pi+m+|C_L\setminus D|.
\]
Hence \(A_L=\pi+m+O(L^{-1})\).

Let \(B\) be any disk with \(|B|=A_L\), and let \(\rho_L:=\sqrt{A_L/\pi}\). Then
\(\rho_L\le R_*\) for some constant \(R_*=R_*(m)\) independent of \(L\). Since
the remote rectangle \(R_L\) lies in the horizontal strip \(|y|\le m/(2L)\), any such disk
intersects the remote rectangle \(R_L\) in area at most
\[
|B\cap R_L|
\le
\frac{2R_*m}{L}.
\]
Also,
\[
|B\cap C_L|\le |C_L|\le \frac{2}{L}.
\]
Therefore
\[
|\Omega_L\cap B|
\le
|D|+|B\cap C_L|+|B\cap R_L|
\le
\pi+\frac{2}{L}+\frac{2R_*m}{L}.
\]
Using \(|B|=|\Omega_L|=A_L\), we get
\[
|\Omega_L\triangle B|
=
2A_L-2|\Omega_L\cap B|
\ge
2m-\frac{C_m}{L}
\]
for a constant \(C_m>0\) depending only on \(m\). Taking the infimum over all
such disks \(B\), we conclude that
\[
\mathcal A(\Omega_L)
\ge
\frac{2m-C_mL^{-1}}{\pi+m}.
\]
Hence
\[
\liminf_{L\to\infty}\mathcal A(\Omega_L)\ge \frac{2m}{\pi+m}.
\]

Now choose \(m>0\) so large that
\[
\frac{2m}{\pi+m}>\alpha.
\]
Then choose \(L\) sufficiently large so that
\[
\|F_L-e_n\|_{\mathcal F^2}<\varepsilon
\qquad\text{and}\qquad
\mathcal A(\Omega_L)\ge \alpha.
\]
Taking \(\Omega:=\Omega_L\) and \(F_\Omega:=F_L\) proves the proposition.
\end{proof}

\begin{remark}
The same argument works for any nonzero scalar multiple of \(w^n\), after
renormalization. In particular, one may make the Fraenkel asymmetry arbitrarily
close to \(2\) while keeping some eigenfunction arbitrarily close, in
\(\mathcal F^2\)-norm, to any prescribed monomial.
\end{remark}

\begin{remark}[Notes toward stability for the perturbative inverse theorem]
The previous proposition shows that a global stability statement for
Theorem~1 in terms of the usual Fraenkel asymmetry cannot hold without further
restrictions: one may keep an eigenfunction arbitrarily close to the reference
monomial while appending very thin remote pieces and corridors, thereby forcing
the ordinary asymmetry to stay large.

There are, however, two natural alternatives that seem compatible with the
perturbative construction.
\end{remark}

\begin{remark}[Weighted local stability]
Fix \(\lambda\in(0,1)\) and let \(D_\lambda\) be the centered Daubechies disk,
that is,
\[
\int_{D_\lambda} e^{-\pi|z|^2}e^{\pi\overline zw}\,dA(z)=\lambda,
\qquad w\in\C.
\]
Let
\[
P(z)=1+\sum_{j=1}^N b_j z^j
\]
with \(\max_j |b_j|\) small, and suppose \(U\) is a bounded \(C^2\) domain
lying in the perturbative class
\[
U_1\subset U\subset U_2
\]
from the proof of Theorem~1, and satisfying
\[
\int_U P(z)e^{-\pi|z|^2}e^{\pi\overline zw}\,dA(z)=\lambda P(w),
\qquad w\in\C.
\]
Then the two implicit-function arguments used in the proof of Theorem~1 should
already yield a local stability statement:
\[
\operatorname{dist}_{C^2}(U,D_\lambda)\lesssim \max_j |b_j|.
\]
Indeed:
\begin{enumerate}
\item the Isakov step produces a unique nearby auxiliary domain \(U_0(P)\),
with \(C^2\) dependence on the datum \(|P|^2\);
\item after applying \(P(\partial_w/\pi)\), the eigenvalue identity becomes the
weighted quadrature identity;
\item the Cherednichenko step produces the unique nearby domain \(U(P)\) having
that weighted Cauchy transform, again with \(C^1\) dependence on the data.
\end{enumerate}
Hence any perturbative solution \(U\) must coincide with that unique branch,
and therefore inherits the same \(C^2\) control. Once \(U\) is a small normal
graph over \(\partial D_\lambda\), one obtains
\[
\int_{U\triangle D_\lambda} e^{-\pi|z|^2}\,dA
\lesssim
\max_j |b_j|.
\]
This weighted asymmetry avoids the remote-rectangle obstruction, since mass
placed far away is exponentially suppressed.

Assessment: this is the most viable part of the program. If the proof of
Theorem~1 already gives a genuinely quantitative implicit-function theorem near
\(D_\lambda\), then the same linearized operator should provide a Lipschitz
estimate for the normal graph parametrization of \(U\) in terms of the trace
data generated by \(P\).

The point that still needs to be written down carefully is the linearization.
One should parametrize nearby boundaries by
\[
\partial U_\rho=\{(R_\lambda+\rho(\theta))e^{i\theta}:\theta\in\mathbb S^1\}
\]
and compute the first variation of the weighted Cauchy data at \(\rho=0\). At
the disk the Fourier modes should diagonalize the derivative, with the first
mode corresponding to translations and the zeroth mode to the change of
\(\lambda\). Once this derivative is shown to be invertible on the orthogonal
complement of those symmetries, the weighted stability estimate should follow
by a standard quantitative implicit-function theorem.
\end{remark}

\begin{remark}[Convex perturbative stability]
If one restricts the class of admissible sets to convex domains, then the same
obstruction is also excluded geometrically. In the perturbative regime just
described, the weighted estimate should imply the ordinary symmetric-difference
estimate
\[
|U\triangle D_\lambda|\lesssim \max_j |b_j|,
\]
because \(U\triangle D_\lambda\) is then confined to a fixed bounded annulus
around \(D_\lambda\), on which the Gaussian weight is bounded below.

Assessment: this perturbative convex statement is also plausible, but only after
the previous \(C^2\) stability has been proved. In that regime the required
annulus confinement is automatic: if \(U\) is a small normal graph over
\(\partial D_\lambda\), then \(U\triangle D_\lambda\subset
\{R_\lambda-c\delta\le |z|\le R_\lambda+c\delta\}\), so the Gaussian weight has
a positive lower bound there.

The main caution is that convexity by itself is not yet doing the work here.
The step from weighted to unweighted asymmetry is immediate once one already
knows that \(U\) stays in a fixed bounded neighborhood of the disk. Thus the
real issue is still the perturbative graph representation, not a separate
convexity estimate.

Thus a reasonable stability statement for Theorem~1 is:
\begin{quote}
near the Gaussian branch, the prescribed eigenfunction determines the domain in
a quantitatively stable way, either in Gaussian-weighted asymmetry for general
perturbative domains, or in ordinary asymmetry on convex perturbative classes.
\end{quote}
\end{remark}

\begin{remark}[What remains to be proved]
To turn the previous notes into a theorem, the main point is to formulate the
perturbative class precisely and then state the uniqueness with quantitative
dependence directly in the language of Theorem~1. If one wants a stronger
convex statement that is not explicitly perturbative, one still needs a
geometric lemma showing that, among convex sets, small weighted asymmetry with
respect to \(D_\lambda\) forces confinement to a bounded annulus around
\(D_\lambda\).

This last step is currently the weakest part of the proposed global statement.
It is not clear from the present notes whether small weighted asymmetry alone
forces such confinement, even under convexity, without some additional control
on the support function or on the area excess.
\end{remark}

\begin{remark}[Program for a global convex stability result]
Suppose now that \(U\) is convex and that one eigenfunction of
\(\mathcal L_U\) is close to a Hermite function \(h_k\). A plausible route to a
global convex stability theorem is the following:
\begin{enumerate}
\item first replace that eigenfunction by a finite Hermite sum, at the price of
slightly perturbing the domain within the convex class;
\item then apply the exact polynomial reduction from the proof of Theorem~1 to
obtain a genuine free-boundary problem;
\item finally use the second proof of the Abreu--D\"orfler theorem to derive a
perturbed holomorphic boundary identity, and convert that identity into
quantitative closeness of \(U\) to a disk by convexity.
\end{enumerate}
This separates the analytic difficulty from the geometric one: once the finite
Hermite reduction has been achieved, the rest of the argument should closely
follow the structure already present in the paper.

Assessment: this is an interesting roadmap, but at present it is still
speculative. The first step is not a soft approximation argument: replacing an
exact eigenfunction by a finite Hermite sum while keeping an exact eigenvalue
equation on a nearby convex domain is essentially a nonlinear realization
problem for the domain-to-eigenfunction map.

Interesting idea: instead of seeking a global finite-mode reduction first, one
could linearize the inverse map directly on the convex side using the support
function \(h(\theta)\) or the radial function \(r(\theta)\). If the boundary
identity from the Abreu--D\"orfler argument can be written as an operator on
\(h\), then one may be able to obtain stability by Fourier analysis on
\(\mathbb S^1\) without introducing an auxiliary exact polynomial eigenfunction.
\end{remark}

\begin{remark}[Candidate finite-mode approximation principle]
The key missing ingredient seems to be a statement of the following type.

\smallskip
\emph{Approximation principle.}
Let \(U\subset\C\) be a bounded convex domain, and let \(F\in\mathcal F^2(\C)\)
be an eigenfunction of \(\mathcal L_U\) with
\[
\|F-e_k\|_{\mathcal F^2}\leq \varepsilon.
\]
Then, after choosing \(N\) large and \(\varepsilon\) small, one should be able
to find a bounded convex domain \(U_N\) and a polynomial
\[
P_N(z)=\sum_{j=0}^N a_j z^j
\]
such that
\[
\mathcal L_{U_N}P_N=\lambda_N P_N,
\]
\[
\|P_N-e_k\|_{\mathcal F^2}\leq C(\varepsilon+\tau_N),
\]
where \(\tau_N\) is the tail of \(F\), and
\[
d(U_N,U)\leq C(\varepsilon+\tau_N)
\]
for a suitable geometric distance \(d\).

\smallskip
The reason this is plausible is that the weighted quadrature-domain step in the
proof of Theorem~1 is already an implicit-function theorem for domains driven by
analytic trace data. If one can compare the trace data attached to \(F\) and to
its truncation \(P_N\), then one may hope to deform \(U\) to a nearby domain
\(U_N\) for which \(P_N\) is an exact eigenfunction. Convexity should then be
preserved by working in a topology where small perturbations keep the domain
convex; in a first approach, this is much more realistic for smooth uniformly
convex domains.

Major problem: nothing in the current notes shows that the image of the
domain-to-trace map contains the truncated data of \(P_N\), even approximately,
in a way compatible with exact solvability. This is the main place where the
argument could fail. In particular, truncating an eigenfunction changes the
entire infinite moment sequence, and it is not obvious that these modified data
still lie on the nonlinear manifold generated by convex domains.
\end{remark}

\begin{remark}[Exact free-boundary reduction after truncation]
Assume that such a pair \((U_N,P_N)\) has been constructed, with \(P_N\) close
to \(z^k\). Then the reduction used in the proof of Theorem~1 applies exactly:
after applying the differential operator \(P_N(\partial_w/\pi)\) to the
eigenvalue identity, one obtains an exact Poisson problem of the form
\[
\Delta u_N
=
-|P_N(z)|^2e^{-\pi|z|^2}\mathbf 1_{U_N}
+\mu_N,
\]
where \(\mu_N\) is an explicit compactly supported distribution depending on the
right-hand side.

This should be viewed as a perturbation of the Hermite free-boundary equation
\[
\Delta u
=
-|z|^{2k}e^{-\pi|z|^2}\mathbf 1_U+c_k\delta_0.
\]
Thus the second proof of the Abreu--D\"orfler theorem suggests that one can
build a holomorphic auxiliary function \(W_N\) in \(U_N\) such that
\[
W_N(z)=c_N-\pi\Phi_k(|z|^2)+E_N(z),
\qquad z\in\partial U_N,
\]
where
\[
\Phi_k(t):=\int_0^t s^k e^{-\pi s}\,ds
\]
and the error \(E_N\) is controlled by the distance between \(P_N\) and
\(z^k\).
\end{remark}

\begin{remark}[What one should extract from the boundary identity]
In the exact Hermite case the error term vanishes, so \(W_N\) is real-valued on
\(\partial U_N\), hence constant in \(U_N\). In the perturbed case one would
hope to prove
\[
\|E_N\|_{L^\infty(\partial U_N)}\leq \eta.
\]
Then the imaginary part of \(W_N\) is \(O(\eta)\) on the boundary, and by the
maximum principle also in the whole domain. Since \(W_N\) is holomorphic, this
should force it to remain close to a real constant, and therefore
\[
\Phi_k(|z|^2)\approx \text{constant}
\qquad\text{for } z\in\partial U_N.
\]
Because \(\Phi_k\) is strictly increasing, this should imply that \(|z|\) is
almost constant on the boundary. At that point the problem becomes geometric.

Major problem: small imaginary part on the boundary does \emph{not} by itself
force a holomorphic function to be close to a constant in any strong norm. One
also needs control of a real normalization and some regularity of the boundary
trace, since the harmonic conjugation operator is not bounded on \(L^\infty\).
So the missing estimate is not merely a maximum-principle step; one needs a
quantitative boundary-regularity argument, probably in \(C^\alpha\),
\(H^{1/2}\), or a similar scale on which the Hilbert transform is controlled.
\end{remark}

\begin{remark}[Convex geometric lemma needed at the end]
The final step should be an elementary convex-geometric statement.

\smallskip
\emph{Radial pinching lemma.}
Let \(U\subset\C\) be a bounded convex domain containing the origin. Assume that
for some \(R>0\) and \(\eta>0\),
\[
\sup_{z\in\partial U}\big||z|-R\big|\leq \eta.
\]
Then
\[
D_{R-\eta}(0)\subset U\subset D_{R+\eta}(0),
\]
and therefore
\[
|U\triangle D_R(0)|\leq C R\,\eta.
\]

\smallskip
Indeed, convexity plus \(0\in U\) implies that every radial segment joining the
origin to a boundary point lies in \(U\), so control of the boundary radii
immediately yields inclusion between concentric disks. Thus, once the holomorphic
argument shows that \(|z|\) is nearly constant on \(\partial U_N\), convexity
upgrades this to quantitative closeness to a disk.

Assessment: this geometric lemma is sound and essentially elementary. The real
content is earlier, namely obtaining a usable estimate of
\(\sup_{\partial U_N}||z|-R|\). Once that is available, convexity converts it
immediately into Hausdorff and symmetric-difference control.
\end{remark}

\begin{remark}[Where the main difficulty lies]
The final geometric step is not the real bottleneck. The difficult point is the
finite-mode approximation principle: one needs a mechanism that replaces a
general eigenfunction \(F\) close to \(e_k\) by an exact finite Hermite
eigenfunction of a nearby convex domain.

The most likely place to search for this is in the proof of the weighted
quadrature-domain preliminary result, since that argument already produces
domains from analytic trace data by an implicit-function theorem. If that
machinery can be re-centered at the actual convex domain \(U\), rather than at a
reference disk, then the truncation \(F\mapsto P_N\) may become compatible with
a controlled deformation \(U\mapsto U_N\).

Interesting idea: recenter the implicit-function theorem at a convex domain
\(U\) by expressing the first variation of the weighted trace map through the
weighted Bergman or Cauchy kernel of \(U\). If one can show that the linearized
map is Fredholm with a controlled inverse in a convexity-preserving topology,
then the finite-mode reduction may become a local projection argument rather
than a global existence problem.
\end{remark}

\section{A Boundary-Oscillation Approach in the Convex Case}

The previous remarks emphasize perturbative implicit-function arguments and
finite-mode reduction. There is, however, a different route that may be better
adapted to convex domains because it keeps the exact eigenvalue equation and
tries to extract rigidity directly from the boundary geometry.

\begin{remark}[Basic setup]
Let \(U\subset\C\) be a bounded convex domain, and suppose that an
\(L^2\)-normalized eigenfunction \(F\) of \(\mathcal L_U\) satisfies
\[
\|F-e_k\|_{\mathcal F^2}\le \varepsilon.
\]
The idea is to parametrize \(U\) by a convexity-compatible boundary variable
and then rewrite the eigenvalue equation as a boundary integral which is exact
for arbitrary convex domains.

If \(0\in U\), one natural parametrization is the radial function
\[
\partial U=\{r(\theta)e^{i\theta}:\theta\in\mathbb S^1\}.
\]
If one wants convexity to enter more linearly, it may be preferable to use the
support function
\[
h(\theta)=\sup_{z\in U}\Re(ze^{-i\theta}),
\]
for which convexity is encoded by the condition \(h''+h\ge 0\) in the weak
sense.
\end{remark}

\begin{remark}[What one gets by integrating by parts]
Starting from
\[
\int_U F(z)e^{-\pi|z|^2}e^{\pi\overline zw}\,dA(z)=\lambda F(w),
\qquad w\in\C,
\]
the first direct attempt is to integrate by parts in \(z\). Writing
\[
\partial_z e^{-\pi|z|^2}=-\pi \overline z\,e^{-\pi|z|^2},
\]
Green's formula gives
\[
\int_U F(z)e^{-\pi|z|^2}e^{\pi\overline zw}\,dA(z)
=
-\frac{1}{2\pi i}\int_{\partial U}
F(z)e^{-\pi|z|^2}e^{\pi\overline zw}\,d\overline z
+\frac{1}{\pi}\int_U F'(z)e^{-\pi|z|^2}e^{\pi\overline zw}\,dA(z).
\]
Iterating \(N\) times yields
\[
\int_U F(z)e^{-\pi|z|^2}e^{\pi\overline zw}\,dA(z)
=
\sum_{j=0}^{N-1}
\frac{(-1)^{j+1}}{2i\,\pi^{j+1}}
\int_{\partial U} F^{(j)}(z)e^{-\pi|z|^2}e^{\pi\overline zw}\,d\overline z
+\frac{1}{\pi^N}\int_U F^{(N)}(z)e^{-\pi|z|^2}e^{\pi\overline zw}\,dA(z).
\]
So the boundary terms are completely explicit, but the procedure does not close
for a general entire eigenfunction: one keeps generating interior terms with
higher derivatives of \(F\). It closes only for polynomial \(F\). Thus, if one
wants a global convex result with no truncation, ordinary integration by parts
is not the right final formula.
\end{remark}

\begin{remark}[The exact boundary identity]
There is an exact contour representation obtained by Cauchy--Green. Set
\[
\Phi_w(z):=F(z)e^{-\pi|z|^2}e^{\pi\overline zw}.
\]
Since
\[
\partial_{\overline z}\Phi_w(z)=\pi(w-z)\Phi_w(z),
\]
one obtains, for every \(w\notin\partial U\),
\[
\int_U \Phi_w(z)\,dA(z)
=
\mathbf 1_U(w)\Phi_w(w)
-\frac{1}{2\pi i}\int_{\partial U}\frac{\Phi_w(z)}{z-w}\,dz.
\]
Because \(\Phi_w(w)=F(w)\), the eigenvalue equation becomes
\begin{equation}\label{eq:global-boundary-identity}
\lambda F(w)
=
\mathbf 1_U(w)F(w)
-\frac{1}{2\pi i}\int_{\partial U}
\frac{F(z)e^{-\pi|z|^2}e^{\pi\overline zw}}{z-w}\,dz.
\end{equation}
Equivalently,
\[
\bigl(\lambda-\mathbf 1_U(w)\bigr)F(w)
=
\frac{1}{2\pi i}\int_{\partial U}
\frac{F(z)e^{-\pi|z|^2}e^{\pi\overline zw}}{w-z}\,dz.
\]
If \(w\in\partial U\), the nontangential boundary value is
\[
\Bigl(\lambda-\frac12\Bigr)F(w)
=
\frac{1}{2\pi i}\,\mathrm{p.v.}\!\int_{\partial U}
\frac{F(z)e^{-\pi|z|^2}e^{\pi\overline zw}}{w-z}\,dz.
\]
This is the exact free-boundary equation on \(\partial U\) for an arbitrary
convex domain.
\end{remark}

\begin{remark}[Boundary parametrizations]
If \(0\in U\) and
\[
z(\theta)=r(\theta)e^{i\theta},
\qquad
dz=\bigl(r'(\theta)+ir(\theta)\bigr)e^{i\theta}\,d\theta,
\]
then \eqref{eq:global-boundary-identity} becomes
\[
\bigl(\lambda-\mathbf 1_U(w)\bigr)F(w)
=
\frac{1}{2\pi i}\int_0^{2\pi}
\frac{
F(r(\theta)e^{i\theta})e^{-\pi r(\theta)^2}e^{\pi r(\theta)e^{-i\theta}w}}
{w-r(\theta)e^{i\theta}}
\bigl(r'(\theta)+ir(\theta)\bigr)e^{i\theta}\,d\theta.
\]

In support-function coordinates one writes
\[
z(\alpha)=e^{i\alpha}\bigl(h(\alpha)+ih'(\alpha)\bigr),
\]
for which
\[
dz=i\,e^{i\alpha}\bigl(h(\alpha)+h''(\alpha)\bigr)\,d\alpha.
\]
Hence
\begin{equation}\label{eq:support-global}
\bigl(\lambda-\mathbf 1_U(w)\bigr)F(w)
=
\frac{1}{2\pi}\int_0^{2\pi}
\frac{F(z(\alpha))e^{-\pi|z(\alpha)|^2}e^{\pi\overline{z(\alpha)}w}}
{w-z(\alpha)}\,e^{i\alpha}\bigl(h(\alpha)+h''(\alpha)\bigr)\,d\alpha.
\end{equation}
This is the most useful global formula, because the geometry enters only
through \(z(\alpha)\) and the nonnegative curvature density \(h+h''\).
\end{remark}

\begin{remark}[What convexity gives when one looks along rays]
Take \(w=t e^{i\beta}\), \(t>0\). In \eqref{eq:support-global}, the exponential
factor has real part
\[
\Psi_{\beta,t}(\alpha)
:=
-\pi|z(\alpha)|^2+\pi t\,\Re\!\bigl(\overline{z(\alpha)}e^{i\beta}\bigr).
\]
Now
\[
\Re\!\bigl(\overline{z(\alpha)}e^{i\beta}\bigr)
=
h(\alpha)\cos(\beta-\alpha)+h'(\alpha)\sin(\beta-\alpha).
\]
Set
\[
S_\beta(\alpha):=
h(\alpha)\cos(\beta-\alpha)+h'(\alpha)\sin(\beta-\alpha).
\]
Then
\[
S_\beta'(\alpha)=(h(\alpha)+h''(\alpha))\sin(\beta-\alpha),
\]
so
\[
S_\beta'(\beta)=0,
\qquad
S_\beta''(\beta)=-(h(\beta)+h''(\beta)).
\]
Therefore convexity turns the boundary integral into a Laplace-type integral
with a stationary point at the support point of outer normal \(e^{i\beta}\). If
\(h(\beta)+h''(\beta)>0\), then
\[
S_\beta(\alpha)
=
h(\beta)-\frac{h(\beta)+h''(\beta)}{2}(\alpha-\beta)^2
+O\bigl(|\alpha-\beta|^3\bigr),
\]
and hence
\[
\Psi_{\beta,t}(\alpha)
=
\pi t h(\beta)-\pi|z(\beta)|^2
-\frac{\pi t}{2}\bigl(h(\beta)+h''(\beta)\bigr)(\alpha-\beta)^2
+O\bigl(t|\alpha-\beta|^3\bigr)+O(|\alpha-\beta|).
\]
So a standard Laplace expansion yields
\[
\bigl(\lambda-\mathbf 1_U(te^{i\beta})\bigr)F(te^{i\beta})
\sim
C(\beta)\,e^{\pi t h(\beta)-\pi|z(\beta)|^2}\,t^{-3/2},
\qquad t\to+\infty,
\]
with an explicit coefficient \(C(\beta)\neq 0\) unless the boundary value
\(F(z(\beta))\) vanishes. The exponent is governed exactly by the support
function \(h(\beta)\).
\end{remark}

\begin{remark}[Assessment of the global convex route]
The preceding computation changes the picture substantially.

First, the exact global object is not a vague oscillatory boundary integral but
the singular Cauchy--Green operator \eqref{eq:global-boundary-identity}. Second,
for general convex sets the relevant asymptotics are not Fourier-mode
oscillations of the whole boundary. They are support-point Laplace asymptotics,
and convexity makes them completely explicit through \(h\) and \(h+h''\).

So a genuinely nonperturbative convex program would be:
\begin{enumerate}
\item start from \eqref{eq:global-boundary-identity} or
\eqref{eq:support-global};
\item study \(w=t e^{i\beta}\) for large \(t\) and extract from the left-hand
side the indicator function \(\beta\mapsto h(\beta)\);
\item compare this indicator with the corresponding growth information carried by
\(F\), using that \(F\) is an eigenfunction close to \(e_k\);
\item prove that the only possible indicator compatible with that comparison is
the constant one, hence \(h\equiv R\) and \(U\) is a disk.
\end{enumerate}

This is the truly interesting global convex problem. The difficult missing input
is now very concrete: one needs a theorem converting the spectral hypothesis
\(\mathcal L_UF=\lambda F\) plus the closeness \(F\approx e_k\) into sharp ray
growth information for \(F(t e^{i\beta})\). Without such a theorem, the support
function asymptotics on the boundary side cannot yet be matched to the right
hand side. But once that growth statement is available, the boundary reduction
above already shows exactly how convexity should force circularity.
\end{remark}

\section{Asymptotic Moment Analysis and a Stability Estimate}

The remarks above identify the key missing input for the global convex program as a theorem converting
$\mathcal{L}_UF=\lambda F$ plus $F\approx e_k$ into sharp ray-growth information for $F$.
We now show that this detour through pointwise asymptotics is unnecessary: the same
conclusion can be reached directly by analysing the matrix elements
$(\mathcal{L}_U)_{jk}$ in polar coordinates and applying a Stirling asymptotic.

\begin{remark}[Off-diagonal bound from the eigenvalue equation]
\label{rem:off-diag}
Let $U\subset\mathbb{C}$ be a bounded domain, $F$ an $L^2$-normalized eigenfunction of $\mathcal{L}_U$
with eigenvalue $\lambda$, and $\|F-e_k\|_{\mathcal{F}^2}\le\varepsilon$.
Write $F=\sum_{n\ge0}c_n e_n$, so $\sum_n|c_n-\delta_{nk}|^2\le\varepsilon^2$.

Because $\mathcal{L}_U$ is a positive contraction ($0\le\mathcal{L}_U\le I$ as operators
on $\mathcal{F}^2$), one has $\|\mathcal{L}_U\|\le1$. Therefore
\[
\|\mathcal{L}_Ue_k-\lambda e_k\|_{\mathcal{F}^2}
=\|(\mathcal{L}_U-\lambda I)(e_k-F)\|_{\mathcal{F}^2}
\le (\|\mathcal{L}_U\|+|\lambda|)\,\|e_k-F\|_{\mathcal{F}^2}
\le 2\varepsilon,
\]
since $(\mathcal{L}_U-\lambda I)F=0$ and $0\le\lambda\le1$.
Expanding in the orthonormal basis $\{e_j\}$:
\[
\|\mathcal{L}_Ue_k-\lambda e_k\|_{\mathcal{F}^2}^2
=\sum_{j\ge0}\bigl|(\mathcal{L}_U)_{jk}-\lambda\delta_{jk}\bigr|^2
\]
where $(\mathcal{L}_U)_{jk}:=\langle\mathcal{L}_Ue_k,e_j\rangle_{\mathcal{F}^2}
=\int_U e_k(z)\overline{e_j(z)}e^{-\pi|z|^2}\,dA(z)$.
In particular,
\begin{equation}\label{eq:off-diag-sum}
\sum_{j\neq k}\bigl|(\mathcal{L}_U)_{jk}\bigr|^2\le4\varepsilon^2.
\end{equation}
\end{remark}

\begin{remark}[Polar-coordinate form of the off-diagonal conditions]
\label{rem:polar-moments}
Suppose $0\in U$ and $U$ is star-shaped, so $\partial U=\{r(\theta)e^{i\theta}:\theta\in[0,2\pi)\}$.
In polar coordinates $z=se^{i\theta}$, $dA=s\,ds\,d\theta$:
\[
(\mathcal{L}_U)_{jk}
=\sqrt{\frac{\pi^{j+k}}{j!\,k!}}\int_0^{2\pi}\!\!e^{i(k-j)\theta}
\!\!\int_0^{r(\theta)}\!\!s^{k+j+1}e^{-\pi s^2}\,ds\,d\theta.
\]
Setting $m:=k-j$ (so $m\neq0$ when $j\neq k$) and defining
\[
G_{k,m}(\theta):=\int_0^{r(\theta)}s^{2k-m+1}e^{-\pi s^2}\,ds,
\]
the off-diagonal element becomes
\[
(\mathcal{L}_U)_{jk}=\sqrt{\frac{\pi^{2k-m}}{(k-m)!\,k!}}\,\widehat{G_{k,m}}(m),
\]
where $\widehat{G_{k,m}}(m):=\int_0^{2\pi}e^{im\theta}G_{k,m}(\theta)\,d\theta$
denotes the $m$-th Fourier coefficient of $G_{k,m}$.

Thus the bound \eqref{eq:off-diag-sum} reads
\begin{equation}\label{eq:fourier-sum}
\sum_{m\neq0}\frac{\pi^{2k-m}}{(k-m)!\,k!}\,\bigl|\widehat{G_{k,m}}(m)\bigr|^2\le4\varepsilon^2,
\end{equation}
an infinite family of weighted Fourier conditions on the boundary function $r(\theta)$.
\end{remark}

\begin{remark}[Stirling asymptotics of the weights]
\label{rem:stirling}
We evaluate the weight $w_{k,m}:=\pi^{2k-m}/((k-m)!\,k!)$ on the circle
$r(\theta)\equiv R:=\sqrt{k/\pi}$, i.e.\ at the disk whose $k$-th eigenfunction
concentrates. Write $R^2=k/\pi$, so $e^{-\pi R^2}=e^{-k}$.

The derivative of $G_{k,m}$ at $r=R$ is
\[
G_{k,m}'(R)=R^{2k-m+1}e^{-\pi R^2}=\bigl(k/\pi\bigr)^{(2k-m+1)/2}e^{-k}.
\]
Hence the weight evaluated on the tangent approximation satisfies
\begin{align*}
w_{k,m}\bigl[G_{k,m}'(R)\bigr]^2
&=\frac{\pi^{2k-m}}{(k-m)!\,k!}
\cdot\frac{k^{2k-m+1}}{\pi^{2k-m+1}}\cdot e^{-2k}
=\frac{k^{2k-m+1}\,e^{-2k}}{\pi\,(k-m)!\,k!}.
\end{align*}
By Stirling's formula, for every fixed $m$ as $k\to\infty$:
\[
(k-m)!\,k!\;\sim\;2\pi\sqrt{k(k-m)}\,\left(\frac{k-m}{e}\right)^{k-m}
\!\!\left(\frac{k}{e}\right)^k
\;\sim\;k^{2k-m}e^{-2k}\quad(\text{leading exponential factor}),
\]
and therefore
\begin{equation}\label{eq:stirling-weight}
w_{k,m}\bigl[G_{k,m}'(R)\bigr]^2\;\sim\;\frac{k}{\pi}
\qquad\text{as }k\to\infty,\;\text{for each fixed }m.
\end{equation}
The same calculation applies for negative $m$ (replacing $k-m$ by $k+|m|$),
with the same leading term $k/\pi$.

Thus every mode $m$ carries weight approximately $k/\pi$ in the quadratic form
\eqref{eq:fourier-sum}, and the weight is \emph{uniform} in $m$ to leading order.
\end{remark}

\begin{remark}[Linearization and the $L^2$ stability estimate]
\label{rem:L2-stab}
Write $r(\theta)=R+\rho(\theta)$ where $R=\sqrt{k/\pi}$ is the reference radius.
For $|\rho(\theta)|\ll R$, expand:
\[
G_{k,m}(\theta)=G_{k,m}(R)+G_{k,m}'(R)\,\rho(\theta)+O\!\bigl(\rho(\theta)^2\,|G_{k,m}''(R)|\bigr).
\]
Since $G_{k,m}'(R)=R^{2k-m+1}e^{-k}$ and
\[
G_{k,m}''(R)=R^{2k-m}e^{-k}\bigl(2k-m+1-2\pi R^2\bigr)=R^{2k-m}e^{-k}(1-m),
\]
the ratio $|G_{k,m}''(R)|/G_{k,m}'(R)=|1-m|/R$ is $O(1/R)$ for fixed $m$.
Hence for $|\rho|\ll R$ the remainder is negligible, and for $m\neq0$:
\[
\widehat{G_{k,m}}(m)=G_{k,m}'(R)\,\hat\rho(m)+O\bigl(\|\rho\|_{L^2}^2/R\bigr),
\]
since the constant term $G_{k,m}(R)$ drops out of the $m\neq0$ Fourier coefficient.

Substituting into \eqref{eq:fourier-sum} and using \eqref{eq:stirling-weight}:
\[
\frac{k}{\pi}\!\sum_{m\neq0}|\hat\rho(m)|^2
\lesssim
\sum_{m\neq0}w_{k,m}\bigl[G_{k,m}'(R)\bigr]^2|\hat\rho(m)|^2
\le4\varepsilon^2.
\]
Identifying the left-hand side as $(k/\pi)\|\rho\|_{L^2(\mathbb{S}^1)}^2$ yields
\begin{equation}\label{eq:L2-bound}
\boxed{\|\rho\|_{L^2(\mathbb{S}^1)}\lesssim\varepsilon\sqrt{\frac{\pi}{k}}.}
\end{equation}

\smallskip
\noindent\emph{Comparison and self-consistency.}
Since $R=\sqrt{k/\pi}$, bound \eqref{eq:L2-bound} reads $\|\rho\|_{L^2}/R\lesssim\varepsilon/k\to0$
for any fixed $\varepsilon$ as $k\to\infty$, confirming that the linearization is self-consistent.
For fixed $k$ the bound gives $\|\rho\|_{L^2}\to0$ as $\varepsilon\to0$.

\smallskip
\noindent\emph{Symmetric-difference control.}
By Cauchy--Schwarz and the inclusion $U\triangle D_R(0)\subset\{|r(\theta)-R|>0\}$:
\begin{align*}
|U\triangle D_R(0)|
&\approx\int_0^{2\pi}R\,|\rho(\theta)|\,d\theta
\le R\sqrt{2\pi}\,\|\rho\|_{L^2}
\lesssim \sqrt{k/\pi}\cdot\sqrt{2\pi}\cdot\varepsilon\sqrt{\pi/k}
=\varepsilon\pi\sqrt{2}.
\end{align*}
Hence $|U\triangle D_R(0)|\lesssim\varepsilon$ with an explicit constant $\pi\sqrt{2}$.
\end{remark}

\begin{remark}[What this computation achieves and what remains]
The preceding calculation establishes a quantitative stability estimate for the
following linearized version of the convex inverse problem:
\begin{quote}
If $U$ is a star-shaped domain and $\|\mathcal{L}_Ue_k-\lambda e_k\|_{\mathcal{F}^2}\le\varepsilon$,
then $|U\triangle D_{\sqrt{k/\pi}}(0)|\lesssim\varepsilon$.
\end{quote}

Several points are worth noting.

\smallskip
\textbf{(a) No ray-growth theorem needed.}
The argument bypasses the ray-growth difficulty identified in the previous section.
The information about $r(\theta)$ is extracted directly from the off-diagonal column
of $\mathcal{L}_U$ in the orthonormal basis, not from the large-$|w|$ behaviour of $F$.

\smallskip
\textbf{(b) The estimate is sharp in $k$.}
The weight $k/\pi$ in \eqref{eq:stirling-weight} is the leading-order coefficient
of the quadratic form, and no cancellation occurs between modes. Thus $\|\rho\|_{L^2}
\sim\varepsilon/\sqrt{k}$ is sharp for this linearized model.

\smallskip
\textbf{(c) What the linearization assumes.}
The step from \eqref{eq:fourier-sum} to \eqref{eq:L2-bound} uses the first-order
Taylor expansion of $G_{k,m}$ at $R$. This requires $|\rho|\ll R$ globally on
$[0,2\pi)$, which is not known a priori from the eigenvalue hypothesis alone.
To remove this assumption one needs either:
\begin{enumerate}
\item a global bound $\|\rho\|_{L^\infty}\le c(k)$ from convexity plus area
normalization (e.g.\ if $|U|=|D_R|$, then $R$ is the correct reference radius
and $|\rho|$ is controlled on average by $\|\rho\|_{L^2}$), or
\item a bootstrapping argument: use \eqref{eq:L2-bound} as a first estimate,
substitute it back to control the nonlinear correction in the Taylor expansion,
and iterate.
\end{enumerate}
For uniformly convex domains of fixed area with a priori bounded curvature, the
first route is immediate: convexity forces $r(\theta)$ to stay within a
bounded annulus around $R$, so the linearization is valid for small $\varepsilon$.

\smallskip
\textbf{(d) The full eigenfunction hypothesis.}
We only used $\|\mathcal{L}_Ue_k-\lambda e_k\|_{\mathcal{F}^2}\le\varepsilon$, which follows
from the eigenfunction hypothesis $\|F-e_k\|_{\mathcal{F}^2}\le\varepsilon$ together
with $\|\mathcal{L}_U\|\le1$.  Thus the same bound
\eqref{eq:L2-bound} holds whenever a normalized eigenfunction of $\mathcal{L}_U$
is $\varepsilon$-close to $e_k$ in $\mathcal{F}^2$, without any further conditions on $\lambda$
or on the eigenvalue gap.

\smallskip
\textbf{(e) Large-$k$ interpretation via incomplete gamma functions.}
One can make the Stirling step more explicit.  The function
$G_{k,0}(\theta)=\int_0^{r(\theta)}s^{2k+1}e^{-\pi s^2}\,ds$
is a re-scaling of the lower incomplete gamma function:
\[
G_{k,0}(\theta)=\frac{1}{2\pi^{k+1}}\,\gamma\!\bigl(k+1,\pi r(\theta)^2\bigr).
\]
The CLT approximation to the Poisson distribution gives, uniformly for
$\pi r^2$ in the range $[k-k^{2/3},k+k^{2/3}]$:
\[
\frac{\gamma(k+1,\pi r^2)}{k!}=\Phi\!\left(\frac{\pi r^2-k}{\sqrt{k}}\right)+O(k^{-1/3}),
\]
where $\Phi$ is the standard normal CDF.  Thus $G_{k,0}(\theta)/G_{k,0}(\infty)$
transitions sharply from $0$ to $1$ on a scale $\Delta r\sim(2\pi)^{-1/2}$ centered
at $r=\sqrt{k/\pi}$.  The derivative at $R=\sqrt{k/\pi}$:
\[
G_{k,0}'(R)=R^{2k+1}e^{-k}
\;\approx\;
\frac{(k!)^{1/2}}{\pi^{k+1/2}}
\cdot\sqrt{\frac{2}{\sqrt{k}}}\cdot\frac{1}{\sqrt{2\pi}}
\]
(using Stirling), which is the same normalization implicit in \eqref{eq:stirling-weight}.
This confirms that \eqref{eq:L2-bound} is an $L^2$ statement about the
level-set function of $r(\theta)$ at the threshold $R=\sqrt{k/\pi}$: the condition
$\sum_m|\hat\rho(m)|^2\lesssim\varepsilon^2\pi/k$ is equivalent, via the co-area formula,
to saying that the $(k+1)$-st Poisson level set of $U$ is within $O(\varepsilon/\sqrt{k})$
in $L^2$ of the circle $\partial D_R$.
\end{remark}

\begin{remark}[Second-order term and its role]
For completeness we record the second-order correction to \eqref{eq:L2-bound}.
At $r=R$ the second derivative $G_{k,m}''(R)=R^{2k-m}e^{-k}(1-m)$ vanishes when
$m=1$.  Thus the mode $m=1$ (the ``translation'' mode) receives only the
third-order contribution from the Taylor expansion.  This is consistent with the
fact that translations of a disk preserve the eigenvalue structure to first order:
translating $D_R$ by $\delta e^{i\beta}$ perturbs $r(\theta)$ by $\delta\cos(\theta-\beta)$,
which is a pure $m=\pm1$ mode, and the first-order change in $\lambda_k$ vanishes
by symmetry.

For $m\neq1$, the second-order term contributes $O(\|\rho\|_{L^2}^4)$ to the
quadratic form in \eqref{eq:fourier-sum}, which is of higher order than the
leading $O(\|\rho\|_{L^2}^2)$ term.  Hence the bound \eqref{eq:L2-bound}
is not affected by the second-order correction in the regime $\|\rho\|_{L^2}\to0$.
\end{remark}

\begin{remark}[Summary of the convex stability program]
Combining the results of this section with the earlier boundary-oscillation analysis,
the following picture emerges.

\emph{Step 1 (Spectral to $L^2$ boundary).}
From $\|F-e_k\|_{\mathcal{F}^2}\le\varepsilon$ and $\mathcal{L}_UF=\lambda F$, one extracts
via Remarks~\ref{rem:off-diag}--\ref{rem:L2-stab}:
\[
\|\,r(\theta)-\sqrt{k/\pi}\,\|_{L^2(\mathbb{S}^1)}\lesssim\varepsilon\sqrt{\pi/k},
\qquad|U\triangle D_{\sqrt{k/\pi}}|\lesssim\varepsilon.
\]

\emph{Step 2 (From $L^2$ to uniform via convexity).}
For uniformly convex $U$, $L^2$ closeness of $r(\theta)$ to a constant $R$ implies
pointwise closeness (since $r$ is Lipschitz with a Lipschitz constant controlled by
the minimum curvature).  The Radial Pinching Lemma of the previous section then
gives $D_{R-\eta}\subset U\subset D_{R+\eta}$ with $\eta\lesssim\varepsilon\sqrt{\pi/k}$,
which implies $|U\triangle D_R|\le C R\eta\lesssim\varepsilon\sqrt{\pi}$.

\emph{What remains.}
For general (not necessarily uniformly convex) convex domains, the passage from
$L^2$ to $L^\infty$ of $r(\theta)$ requires controlling the modulus of convexity.
A natural condition sufficient for this is: $U$ has area equal to $|D_R|=k/\pi$
and support function satisfying $h''+h\ge c>0$ uniformly.
Under this curvature bound, the Fourier decay of $r(\theta)$ is $O(1/m^2)$,
and one can close the argument without further hypotheses.
\end{remark}

\section{A Compactness Route to Global Convex Stability}

\begin{remark}[Why the previous argument is only local]
The estimate in Remark~\ref{rem:L2-stab} is fundamentally perturbative: it gives
strong control once one already knows that the radial function $r(\theta)$ stays
in a small neighborhood of the reference radius $R=\sqrt{k/\pi}$. What it does
\emph{not} by itself rule out is the existence of a sequence of convex sets
$U_N$ of the correct area for which an eigenfunction of $\mathcal L_{U_N}$ is
very close to $e_k$, while the geometry of $U_N$ drifts far from the disk, for
instance by becoming thinner and more elongated.

Thus a full global theorem should be split into two parts:
\begin{enumerate}
\item a \emph{compactness/non-degeneration step}, whose purpose is to prevent
convex sets with almost-monomial eigenfunctions from escaping to infinity in the
moduli space of convex bodies, and
\item a \emph{local quantitative step}, namely the perturbative analysis above,
which converts qualitative closeness to the disk into an explicit estimate.
\end{enumerate}
The point is that the present section essentially settles the second part once
the first one is available.
\end{remark}

\begin{remark}[A concrete contradiction scheme]
Fix $k\in\N\cup\{0\}$ and set $R=\sqrt{k/\pi}$, so that $|D_R|=k$.
Suppose the desired global stability statement were false. Then there would exist
an $\eta_0>0$ and a sequence of bounded convex sets $U_N\subset\C$ with
\[
|U_N|=|D_R|,
\qquad
\mathcal A(U_N)\ge \eta_0,
\]
and normalized eigenfunctions $F_N$ of $\mathcal L_{U_N}$ such that
\[
\|F_N-e_k\|_{\mathcal F^2}\to 0.
\]
After multiplying by unimodular constants one may assume the phase is fixed.

To reach a contradiction one would like to make the following five steps
completely explicit.

\smallskip
\textbf{Step 1: choose a normalization of the sets.}
Since translations act nontrivially on monomials, one must pick a canonical
center for $U_N$ before taking limits. There are several natural choices:
the barycenter of $U_N$, the center of the best-fitting disk in the Fraenkel
asymmetry, or the center that kills the $m=\pm1$ Fourier mode of the boundary.
What matters is that the normalization be stable under convergence and compatible
with the fact that the target eigenfunction is the \emph{centered} monomial $e_k$.

\smallskip
\textbf{Step 2: prove a non-degeneration estimate for convex bodies.}
One needs an a priori statement saying that a convex set with an eigenfunction
very close to $e_k$ cannot be arbitrarily elongated. Concretely, after the
normalization of Step~1, one wants a bound of the form
\[
\operatorname{diam}(U)\le C(k,\delta)
\]
whenever $U$ is convex, $|U|=|D_R|$, and $\mathcal L_U$ has a normalized
eigenfunction $F$ with $\|F-e_k\|_{\mathcal F^2}\le \delta$.
Equivalently, if $\operatorname{diam}(U_N)\to\infty$ while $|U_N|=|D_R|$, then
every normalized eigenfunction of $\mathcal L_{U_N}$ must stay a definite
distance away from $e_k$ in $\mathcal F^2$. This is precisely the place where
one would show that ``long and thin'' convex sets have eigenfunctions that are
too anisotropic to look Hermite-like.

\smallskip
\textbf{Step 3: compactness of the shapes and continuity of the operators.}
Once Step~2 gives a uniform diameter bound, Blaschke's selection theorem yields,
after passing to a subsequence, Hausdorff convergence
\[
U_N\to U_\infty
\]
to a bounded convex set of the same area. In particular,
\[
|U_N\triangle U_\infty|\to 0.
\]
Since localization operators depend Lipschitz-continuously on the symbol in
area measure,
\[
\|\mathcal L_{U_N}-\mathcal L_{U_\infty}\|_{\mathcal F^2\to\mathcal F^2}
\le |U_N\triangle U_\infty|,
\]
it follows that $\mathcal L_{U_N}\to\mathcal L_{U_\infty}$ in operator norm.

\smallskip
\textbf{Step 4: identify the limiting set.}
Let $\lambda_N$ be the eigenvalue corresponding to $F_N$.
After extraction, $\lambda_N\to\lambda_\infty\in[0,1]$. Using
\[
\mathcal L_{U_N}F_N=\lambda_NF_N,
\qquad
\|F_N-e_k\|_{\mathcal F^2}\to 0,
\qquad
\|\mathcal L_{U_N}-\mathcal L_{U_\infty}\|\to 0,
\]
one obtains
\[
\|\mathcal L_{U_\infty}e_k-\lambda_\infty e_k\|_{\mathcal F^2}\to 0,
\]
hence
\[
\mathcal L_{U_\infty}e_k=\lambda_\infty e_k.
\]
At this point one invokes the exact inverse result proved earlier in the paper:
if a convex set has $e_k$ as an eigenfunction, then that set must be the centered
disk $D_R$. Therefore $U_\infty=D_R$.

\smallskip
\textbf{Step 5: contradict the assumption of non-stability.}
Because $U_N\to D_R$ in Hausdorff (hence also in symmetric-difference) distance,
their asymmetry satisfies
\[
\mathcal A(U_N)\to 0,
\]
contradicting $\mathcal A(U_N)\ge\eta_0$. This proves a \emph{qualitative global
stability statement}: for every $\eta>0$ there exists $\delta(\eta,k)>0$ such
that if $U$ is convex, $|U|=|D_R|$, and $\mathcal L_U$ has a normalized
eigenfunction $F$ with $\|F-e_k\|_{\mathcal F^2}\le\delta(\eta,k)$, then
$\mathcal A(U)<\eta$.
\end{remark}

\begin{remark}[Where the real work sits]
The compactness scheme above is clean; the entire difficulty is concentrated in
Step~2. One needs a mechanism forcing a spectral penalty on elongated convex sets.
A possible target statement is the following:
\begin{quote}
For every $\eta>0$ there exists $c_{\eta,k}>0$ such that, if $U$ is convex,
$|U|=|D_R|$, and $\mathcal A(U)\ge\eta$, then every normalized eigenfunction
$F$ of $\mathcal L_U$ satisfies
\[
\inf_{\gamma\in\R}\|F-e^{i\gamma}e_k\|_{\mathcal F^2}\ge c_{\eta,k}.
\]
\end{quote}
This is the exact ``no almost-monomials far from disks'' statement needed to
start the contradiction argument.

There are several plausible ways to attack it.

\smallskip
\textbf{(a) Width/diameter coercivity from convexity.}
For a convex set of fixed area, large diameter forces a very small width in the
transverse direction. One would then try to show that the moments
\[
\int_U z^j\overline{z}^{\,\ell}e^{-\pi|z|^2}\,dA(z)
\]
cannot mimic the disk moments when the body has such a one-dimensional profile.
In other words, extreme eccentricity should show up in the off-diagonal matrix
elements of $\mathcal L_U$ in the monomial basis.

\smallskip
\textbf{(b) A selection principle in the spirit of Brasco--De Philippis--Velichkov.}
Instead of working directly with an arbitrary bad sequence, one could minimize a
penalized functional of the form
\[
U\mapsto
\frac{\inf\{\|F-e_k\|_{\mathcal F^2}^2:\ \|F\|=1,\ \mathcal L_UF=\lambda F\}}
{\mathcal A(U)^2}
+\Lambda\,\Phi(U),
\]
over convex sets of fixed area, where $\Phi(U)$ is a secondary coercive term
controlling translation or diameter. The Euler--Lagrange information for almost
minimizers may be more rigid than for an arbitrary minimizing sequence, just as
in the sharp Faber--Krahn strategy.

\smallskip
\textbf{(c) Separate compactness from quantification.}
Once the above contradiction argument gives qualitative closeness $U\approx D_R$,
the perturbative computation of Remark~\ref{rem:L2-stab} should become applicable.
One then recovers a quantitative estimate by running the local argument around
the disk. In this sense, the global theorem naturally decomposes as
\[
\text{compactness }+\text{ exact characterization }+\text{ local expansion}.
\]
This is conceptually the same structure as in many sharp spectral stability
results: first exclude escaping sequences, then identify the only possible limit,
and only after that linearize near the optimizer.
\end{remark}

\section{A Formal One-Dimensional Picture for Thin Rectangles}

\begin{remark}[Heuristic reduction for the elongated rectangle]
Consider the area-one rectangles
\[
R_L:=
\left[-\frac L2,\frac L2\right]\times
\left[-\frac{1}{2L},\frac{1}{2L}\right],
\qquad L\to\infty.
\]
The small-norm lemma shows that all eigenvalues of $\mathcal L_{R_L}$ tend to
$0$, so there is no nontrivial limiting eigenfunction at fixed scale.
Nevertheless, one can still ask what the eigenfunctions \emph{look like} before
normalizing away the small eigenvalue.

The natural heuristic is that, because the vertical width is $L^{-1}$, the
operator should become effectively one-dimensional: the relevant modes vary
along the long horizontal direction and are approximately constant across the
thin vertical direction.
\end{remark}

\begin{remark}[Formal derivation of the effective kernel]
Write $z=x+iy$ and $w=u+iv$. Then
\[
(\mathcal L_{R_L}F)(w)
=
\int_{-L/2}^{L/2}\int_{-1/(2L)}^{1/(2L)}
F(x+iy)e^{\pi(x-iy)(u+iv)}e^{-\pi(x^2+y^2)}\,dy\,dx.
\]
Since $|y|\le 1/(2L)$, the thin direction is frozen at leading order. Formally,
for functions varying on $O(1)$ scales one may write
\[
F(x+iy)=F(x)+O(L^{-1}\partial_yF),
\qquad
e^{-\pi y^2}=1+O(L^{-2}),
\qquad
e^{\pi(-iy)(u+iv)}=1+O(L^{-1}|w|).
\]
Substituting this into the integral gives
\[
(\mathcal L_{R_L}F)(w)
=
\frac1L\int_{-L/2}^{L/2}F(x)e^{\pi xw}e^{-\pi x^2}\,dx
+\text{lower-order terms}.
\]
Thus, to first order, the dependence on the thin variable $y$ disappears.

Equivalently, if one embeds one-variable functions by
\[
(J_Lf)(x+iy):=f(x),
\]
then one expects
\[
\mathcal L_{R_L}J_L\approx J_LT_L,
\]
where $T_L$ is an effective one-dimensional integral operator acting on
$[-L/2,L/2]$.
\end{remark}

\begin{remark}[Gaussian conjugation and the effective one-dimensional operator]
After the usual Gaussian conjugation, the effective kernel takes a symmetric
form. More precisely, writing
\[
f(x):=F(x)e^{-\pi x^2/2},
\]
the factor
\[
e^{\pi xu}e^{-\pi x^2}
=
e^{-\frac{\pi}{2}(x-u)^2}e^{\frac{\pi}{2}(u^2-x^2)}
\]
shows that the reduced operator is formally conjugate to
\[
(T_Lf)(u)
:=
\frac1L\int_{-L/2}^{L/2}e^{-\frac{\pi}{2}(u-x)^2}f(x)\,dx.
\]
Hence the asymptotic spectral problem for $\mathcal L_{R_L}$ should be governed,
to leading order, by the compact self-adjoint operator with kernel
\[
K_L(u,x):=\frac1L e^{-\frac{\pi}{2}(u-x)^2}
\mathbf 1_{[-L/2,L/2]}(u)\mathbf 1_{[-L/2,L/2]}(x).
\]

In particular, if $(\mu_{n,L},\phi_{n,L})$ denotes the $n$-th eigenpair of
$T_L$, then the corresponding two-dimensional eigenfunction should satisfy the
formal approximation
\[
F_{n,L}(x+iy)\approx \phi_{n,L}(x),
\]
that is, it is essentially constant across the thin direction and oscillates
only along the long axis.
\end{remark}

\begin{remark}[What this predicts for the eigenfunctions]
The previous reduction suggests the following qualitative picture.

\smallskip
\textbf{(a) Principal mode.}
The first eigenfunction is positive, even in $x$, nearly constant on the bulk of
$[-L/2,L/2]$, and decays only near the endpoints $x=\pm L/2$.

\smallskip
\textbf{(b) Higher modes.}
The higher eigenfunctions are alternating even and odd longitudinal modes, with
successively more sign changes in the horizontal variable $x$, while remaining
almost constant in the transverse variable $y$.

\smallskip
\textbf{(c) Failure of disk-like behavior.}
These modes are strongly anisotropic and therefore are not expected to resemble
the radial or monomial eigenfunctions attached to disks. In particular, very
elongated rectangles should force eigenfunctions to look ``one-dimensional''
rather than Hermite-like in the two-dimensional Bargmann--Fock sense.
\end{remark}

\begin{remark}[Bulk asymptotics for fixed mode number]
For fixed $n$ and $L\to\infty$, the interval $[-L/2,L/2]$ is locally large, so
away from the endpoints the operator $T_L$ is close to full Gaussian convolution
on $\R$:
\[
(T_\infty f)(u)=\frac1L\int_{\R}e^{-\frac{\pi}{2}(u-x)^2}f(x)\,dx.
\]
The latter is diagonalized by Hermite functions after rescaling. Therefore one
expects that, for each fixed $n$ and for $x$ in the bulk region
$|x|\le L/2-o(L)$,
\[
\phi_{n,L}(x)\approx h_n\!\left(\sqrt{\frac{\pi}{2}}\,x\right),
\]
where $h_n$ denotes the $n$-th normalized Hermite function (with the usual
convention on $\R$). Consequently,
\[
F_{n,L}(x+iy)\approx h_n\!\left(\sqrt{\frac{\pi}{2}}\,x\right)
\]
in the bulk, with the endpoint corrections encoded by the truncation to
$[-L/2,L/2]$.

Thus the clean formal summary is
\[
\boxed{
F_{n,L}(x+iy)\approx \phi_{n,L}(x),
\qquad
\mu_{n,L}\phi_{n,L}(u)=\frac1L\int_{-L/2}^{L/2}
e^{-\frac{\pi}{2}(u-x)^2}\phi_{n,L}(x)\,dx,
}
\]
and, for fixed mode number $n$,
\[
\boxed{
F_{n,L}(x+iy)\approx h_n\!\left(\sqrt{\frac{\pi}{2}}\,x\right)
\quad\text{in the bulk.}
}
\]
This should be understood as a formal asymptotic guide, not as a theorem proved
in this note.
\end{remark}

\end{document}
